\documentclass[11pt,a4paper]{article}
\usepackage[margin=2.5cm]{geometry}
\usepackage{amsmath,amssymb}
\usepackage[T1]{fontenc}
\usepackage{lmodern}
\usepackage{microtype}
\usepackage{enumitem}
\setlist{itemsep=2pt,topsep=3pt}
\newcommand{\TTb}{$T\bar T$}

\title{Report on \emph{``Matrix Theory Reloaded: A BPS Road to Holography''}\\[4pt]
\large by C.~D.~A.~Blair, J.~Lahnsteiner, N.~A.~Obers and Z.~Yan}
\date{6 November 2026}

\begin{document}
\maketitle

\section*{Summary of the paper}

The paper treats the decoupling limit that produces matrix theory on D$p$-branes as a
\emph{BPS decoupling limit of the whole of type~II string theory}, not just of the brane.
The resulting corner, ``matrix $p$-brane theory'' (M$p$T), has a non-Lorentzian target
space with a codimension-$(p+1)$ foliation. It is described by $p$-brane Newton--Cartan
vielbeins and has a $p$-brane Galilean boost symmetry. The paper then develops three strands.
\begin{enumerate}[label=(\roman*)]
\item \emph{Holography.} The paper applies the M$p$T limit at asymptotic infinity of a D$q$-brane supergravity
  solution. For $p=q$ this limit \emph{induces} the near-horizon (Maldacena/IMSY) limit in the bulk. For
  $p\neq q$ it gives a non-Lorentzian bulk geometry in which a second, M$q$T-type limit
  produces a near-horizon geometry. This reproduces the recent non-Lorentzian proposals of Lambert
  et al.\ in a systematic way.
\item \emph{DLCQ conjecture.} Lifted to M-theory, each BPS limit corresponds to a DLCQ,
  up to U-duality. This motivates the conjecture that holography is a
  DLCQ$^n$/DLCQ$^m$ correspondence with $m>n$. In this language AdS$_5$/CFT$_4$ is
  DLCQ$^0$/DLCQ$^1$, the non-Lorentzian examples are DLCQ$^1$/DLCQ$^2$, and D1--D5
  AdS$_3$/CFT$_2$ is DLCQ$^0$/DLCQ$^2$.
\item \emph{\TTb.} Undoing the BPS limit is identified with the \TTb{} deformation for the
  string, and with U-dual ``$p$-brane \TTb'' deformations for branes. This yields explicit
  flow equations:
  \begin{itemize}
  \item a closed-form flow in any dimension, taking $D$ free scalars to the Dirac--Nambu--Goto action;
  \item a strikingly simple cubic flow for $p=2$;
  \item flows including the gauge field for $p=1,2$;
  \item perturbative flows for $p\geq 3$.
  \end{itemize}
\end{enumerate}

\section*{1. Principal achievement}

The main achievement is conceptual unification: the \emph{near-horizon limit is not an
independent ingredient of holography}. It is the bulk image of the BPS decoupling limit that defines
the boundary theory, and the near-horizon geometry automatically approaches the flat
non-Lorentzian M$p$T geometry at its own asymptotic boundary. Matrix theory, AdS/CFT and its
IMSY generalisations, and the recently proposed non-Lorentzian holographic dualities were
until now treated with separate, partly ad hoc scalings. Here they become entries of a single
duality-organised table, generated by one rule applied repeatedly. The analysis also explains
\emph{why} only certain brane combinations work: $q=p+4$ for aligned configurations, and
$p+q=10$ in general, unless one smears or rescales $N$. The bookkeeping in terms of iterated
DLCQs gives a common vocabulary for flat-space, non-Lorentzian and (via timelike T-duality)
Carrollian corners.

The most portable concrete result is the family of generalised \TTb{} flow equations,
\[
\frac{\partial\mathcal L}{\partial\lambda}=\frac{1}{2\lambda^2}\Big\{\mathrm{tr}(\mathbb 1-\lambda\mathcal T)
-(p-1)\big[\det(\mathbb 1-\lambda\mathcal T)\big]^{\frac1{p-1}}-2\Big\}.
\]
This equation reproduces the known $p=0$ (Gross et al.) and $p=1$ (standard \TTb) cases. For
$p=2$ it gives a new and simple cubic flow. Unlike earlier higher-dimensional proposals, it
reaches the full multi-scalar Nambu--Goto action, and it explains why ordinary \TTb{} acting
on 2D Maxwell theory does not reach DBI.

I checked the equation numerically. Along the claimed solution family, the closed-form
Nambu--Goto-type Lagrangian, both sides agree to numerical precision for $p=0,2,3,4,5$ with 1,
3 and 6 scalars. The $p=2$ polynomial form agrees as well. The claims are therefore correct as
stated.

\section*{2. Principal drawback}

The principal drawback is that the framework is \emph{kinematic}: it reorganises limits and
geometries, but no new duality or new deformation is put to a \emph{dynamical} test. Several
features of the paper show this.
\begin{itemize}
\item \emph{The established dualities are re-described, not extended.} For $p=q$ the asymptotic
  M$p$T limit is shown (Sec.~3.1) to coincide with the Maldacena limit after a
  reparametrisation.
\item \emph{The new non-Lorentzian bulk backgrounds are not established.}
  \begin{itemize}
  \item They are not shown to solve the equations of motion and torsion constraints of
    M$p$T supergravity, which the authors acknowledge.
  \item Their isometries and supersymmetries are not computed.
  \item Generic cases need smearing, or a by-hand $\omega$-dependent rescaling of $\ell$ (that
    is, of $N$), whose physical meaning is left open.
  \end{itemize}
\item \emph{No observable is matched on both sides.} No spectrum, symmetry algebra,
  thermodynamic quantity, protected correlator or index is computed for any proposed new duality.
  The DLCQ$^n$/DLCQ$^m$ conjecture therefore works as a classification scheme without a
  criterion that could falsify it. Its DLCQ$^2$/DLCQ$^3$ example (Spin Matrix Theory) rests
  on an admittedly tentative identification of the third DLCQ.
\item \emph{The \TTb{} ``intrinsic generation'' of the bulk is partly circular.}
  \begin{itemize}
  \item The radial profile $\Omega(r)$ is inserted by hand using the dilatation symmetry; the
    deformation does not determine it.
  \item The $p$-brane flows are classical, bosonic and abelian.
  \item They are fixed ``by construction'' only along the DBI family. The operator is not
    unique off this family: in the single-scalar case the paper's flow coincides with the
    inequivalent flow of Ferko et al.
  \item There is no quantum definition, and no analogue of Zamolodchikov factorisation.
  \end{itemize}
\end{itemize}
Above all, the step that matters for holography is missing: applying the deformation to the
large-$N$, non-abelian worldvolume theory. As a result, the paper's answer to its own Question~3
is currently a statement about a single classical brane action, not about the strongly coupled
theory dual to the bulk. A secondary issue is length. The paper runs to roughly 100 pages, much
of it review, and this dilutes the genuinely new content.

\section*{3. Most worthwhile follow-up}

\paragraph{Proposal: a large-$N$ thermodynamic test of ``$p$-brane \TTb{} = undoing the
near-horizon limit'', starting with the M2-brane ($p=2$).}

\emph{Why this case.} The $p=2$ flow is exact, polynomial and independent of the field
strength. The paper derives it for both D2- and M2-branes, and M-theory has no running dilaton.
It is therefore the cleanest place to test the claim that the deformation interpolates between
the near-horizon AdS$_4\times S^7$ geometry and the full, asymptotically flat brane solution
(the ``1'' in $H=1+L^6/r^6$ restored).

\emph{Plan.}
\begin{enumerate}[label=(\alph*)]
\item \emph{Field-theory side.} Apply the cubic flow to ABJM on $T^2$ at large $N$. Use
  large-$N$ factorisation of the composite operators
  $\mathrm{tr}\,\mathcal T^2$, $(\mathrm{tr}\,\mathcal T)^2$ and $\det\mathcal T$, following
  the treatment of Taylor's operator by Hartman, Kruthoff, Shaghoulian and Tajdini. This gives
  a closed, Burgers-type equation for $E(\lambda;S,V)$ in homogeneous states, which can be
  integrated from the conformal equation of state $\varepsilon\propto s^{3/2}$.
\item \emph{Gravity side.} Compute the thermodynamics of non-extremal M2-branes in
  asymptotically flat space at fixed charge, with $\lambda$ identified through
  $\omega^{-2}$ and the paper's dictionary. One striking feature makes the test promising. At
  fixed charge, $T=6/(4\pi r_0\cosh\beta)$ is bounded and has two branches: a near-extremal
  branch that reduces to AdS$_4$ thermodynamics, and an unstable far-from-extremal branch.
  This mirrors the maximal (Hagedorn-like) temperature and square-root branch structure of
  \TTb{} with the corresponding sign.
\item \emph{Compare} the free energies $F(T,\lambda)$, including the location of $T_{\max}$.
  Then extend the comparison to:
  \begin{itemize}
  \item the D3-brane, using the scalar flow plus the perturbative $F$-terms;
  \item the D1--D5 system with its double deformation. This is S-dual to the single-trace
    \TTb{} of Giveon--Itzhaki--Kutasov, where an analogous match is known to work.
  \end{itemize}
\end{enumerate}

\emph{Feasibility and impact.} Both sides are classical or large-$N$ computations with
established tools, so the project is realistic on a time-scale of months. Either outcome is
informative.
\begin{itemize}
\item A match would give the first dynamical evidence for the paper's central claim. It would
  also supply a higher-dimensional, \TTb-like deformation with a concrete asymptotically flat
  holographic dual, which matters for flat-space holography.
\item A mismatch would show whether the deformation must be implemented in a ``single-trace''
  rather than a factorised (``double-trace'') way. It would thereby constrain the missing
  non-abelian completion of the flow.
\end{itemize}
The main risks are two. Factorisation of the cubic $\det\mathcal T$ term in $d=3$ needs
justification, and the map between $\lambda$ and the supergravity parameters must be fixed
with care. Neither looks prohibitive. The project addresses the paper's principal weakness
directly, at modest cost.

\end{document}
